\documentclass[11pt]{article}

% ============================================================
% Milestone 1 -- Project Proposal & Literature Review Template
% DIP Team Project
% ============================================================
\usepackage[margin=1in]{geometry}
\usepackage[T1]{fontenc}
\usepackage[utf8]{inputenc}
\usepackage{amsmath,amssymb}
\usepackage{graphicx}
\usepackage{booktabs}
\usepackage[round]{natbib}
\usepackage{hyperref}
\usepackage{float}
\hypersetup{colorlinks=true, linkcolor=blue!50!black, citecolor=blue!50!black, urlcolor=blue!50!black}
\usepackage{enumitem}
\usepackage{xcolor}

% TikZ libraries required for the pipeline diagram
\usepackage{tikz}
\usetikzlibrary{shapes.geometric, arrows.meta, positioning, calc}

\newcommand{\todo}[1]{\textcolor{red}{\textbf{[TODO: #1]}}}

\title{\LARGE Project Proposal \& Literature Review \\[4pt]
\large Compressed-Domain Image Classification}
\author{
  Habiba Mounir \and
  \todo{Student B Full Name} \and
  \ Shahd Ahmed
}
\date{Milestone 1 --- \todo{Date} \\{} [CSE681] --- Digital Image Processing}

\begin{document}
\maketitle

\begin{abstract}
Image classification systems typically decode JPEG images into RGB pixels before inference. This project investigates whether classification can instead be performed directly on JPEG-domain DCT coefficients. We will compare a DCT-domain neural network against a ResNet-50 trained on decoded RGB images using the CIFAR-10 dataset. Evaluation will focus on classification accuracy and computational efficiency in terms of inference latency and preprocessing cost.
\end{abstract}

\section{Problem Statement \& Motivation}
\label{sec:motivation}
\todo{What image-processing problem are you solving? Why is it interesting or useful?
How does it connect to material from the course (name the week/topic)? 1 paragraph.}

\section{Related Work}
\label{sec:related}
\ To build a robust foundation for our compressed-domain image classification approach, our project draws upon recent advancements and established architectural frameworks in efficient deep learning. We primarily rely on the foundational methodology established by Gueguen et al. (2018), which proves the viability of bypassing traditional pixel decompression to classify directly from JPEG DCT coefficients. Furthermore, to address modern hardware constraints and latency benchmarks, we incorporate insights from Canziani et al. (2016) regarding resource profiling on laptop-class hardware.

In addition to foundational works, recent literature from 2023 and 2024—such as the efficient compressed-domain processing techniques by Wu et al. (2023) and the JPEG stream acceleration methods by Zhang et al. (2024)—validate our strategy of trading off high-frequency DCT bands to achieve optimal end-to-end latency without sacrificing classification accuracy. These studies directly inform our ablation study protocol and hardware standardization framework.

\begin{table}[H]
\centering
\caption{Key references.}
\label{tab:refs}
\begin{tabular}{@{}p{0.30\linewidth}p{0.12\linewidth}p{0.45\linewidth}@{}}
\toprule
\textbf{Reference} & \textbf{Year} & \textbf{Relevance to this project} \\
\midrule
Gueguen et al. \citep{gueguen2018faster} & 2018 & Provides the core architecture for direct JPEG-DCT feature classification without decompression. \\ \hline
Canziani et al. \citep{canziani2016analysis} & 2016 & Guides our hardware profiling and end-to-end latency benchmarks on resource-constrained setups. \\ \hline
Wu et al. \citep{wu2023efficient} & 2023 & Offers modern insights into efficient compressed-domain image processing and optimization. \\ \hline
Zhang et al. \citep{zhang2024accelerating} & 2024 & Inspires our frequency band ablation studies and stream acceleration techniques. \\ \hline
\bottomrule
\end{tabular}
\end{table}

\section{Dataset \& Baseline Plan}
\label{sec:data}
The CIFAR-10 dataset will be used. It contains 50,000 training images and 10,000 test images across ten object categories. Images will be re-encoded as JPEG quality 90 before DCT extraction. The baseline model is a standard ResNet-50 operating on decoded RGB images. The proposed model will operate directly on the JPEG DCT coefficients extracted from the same images.

\begin{table}[H]
\centering
\caption{Dataset summary.}
\label{tab:data}
\begin{tabular}{@{}lll@{}}
\toprule
\textbf{Dataset} & \textbf{Subset used} & \textbf{Link} \\
\midrule
CIFAR-10 & 50,000 train, 10,000 test & https://www.cs.toronto.edu/~kriz/cifar.html \\
\bottomrule
\end{tabular}
\end{table}

\begin{figure}[H]
  \centering
  \tikzset{
      block/.style = {rectangle, draw, fill=blue!10, text width=6.8cm, text centered, rounded corners, minimum height=0.9cm, font=\small},
      process/.style = {rectangle, draw, fill=orange!15, text width=6.8cm, text centered, minimum height=0.9cm, font=\small},
      data/.style = {trapezium, trapezium stretches body, fill=green!15, draw, text width=5.8cm, text centered, minimum height=0.8cm, font=\small},
      arrow/.style = {thick, -{Stealth[scale=1.2]}}
  }
  
  \begin{tikzpicture}[node distance=1.8cm]
      % Stage 1: CIFAR-10 to JPEG
      \node (cifar) [data] {CIFAR-10 Dataset\\(50K Train / 10K Test Arrays)};
      
      \node (jpeg_save) [process, below=of cifar] {Convert NumPy Arrays to PIL Images\\\& Save as JPEG (Quality = 100)\\Structured by Class Labels};
      
      \node (jpeg_dir) [data, below=of jpeg_save] {JPEG Dataset Disk Storage\\(\texttt{jpeg\_dataset/train/} \& \texttt{test/})};
      
      % Stage 2: JPEG to DCT .npy
      \node (jio_read) [process, below=of jpeg_dir] {Parse JPEG Structs via \texttt{jpegio}\\Extract Quantized DCT Arrays (\texttt{coef\_arrays})};
      
      \node (dequant) [process, below=of jio_read] {Dequantize Coefficients per Channel:\\\texttt{coef\_dequant = coef\_q * qtbl\_full}\\(Y, Cb, Cr matrices)};
      
      \node (npy_save) [data, below=of dequant] {DCT Compressed-Domain Dataset\\(\texttt{*.npy} Dictionaries: Y, Cb, Cr)};

      % Connectors
      \draw [arrow] (cifar) -- (jpeg_save);
      \draw [arrow] (jpeg_save) -- (jpeg_dir);
      \draw [arrow] (jpeg_dir) -- (jio_read);
      \draw [arrow] (jio_read) -- (dequant);
      \draw [arrow] (dequant) -- (npy_save);
  \end{tikzpicture}
  \caption{Data preprocessing pipeline illustrating the translation of raw CIFAR-10 pixel batches into dequantized DCT coefficients saved as compressed dictionary blocks.}
  \label{fig:pipeline_flowchart}
\end{figure}

\section{Proposed Method}
\label{sec:method}
\todo{Sketch your planned approach: architecture diagram or pseudocode, what is novel
or extended relative to the baseline, and why you expect it to work.}

\section{Feasibility \& Compute Plan}
\label{sec:feasibility}
\todo{State your compute tier (CPU-OK / GPU-light), expected training time per
experiment, and confirm this fits a laptop-class machine within the project window.}

\subsection{Minimum Viable Deliverable (MVP)}
Develop and evaluate a JPEG-domain image classification pipeline on CIFAR-10. The MVP will:

\begin{enumerate}
    \item Re-encode CIFAR-10 images as JPEG.
    \item Extract Y/Cb/Cr DCT coefficients.
    \item Train a ResNet-50 RGB baseline.
    \item Train a DCT-domain classifier.
    \item Measure classification accuracy, latency, and throughput.
\end{enumerate}

Success is defined as achieving at least 88\% top-1 classification accuracy on CIFAR-10 while demonstrating a measurable reduction in inference cost compared to the RGB baseline.

\subsection{Stretch Goal}
\todo{What you will attempt if the MVP is on track by Week 5--6.}

\section{Risk Assessment}
\label{sec:risk}
\begin{itemize}[leftmargin=*]
  \item \todo{Risk 1 and its mitigation / fallback.}
  \item \todo{Risk 2 and its mitigation / fallback.}
\end{itemize}

\section{Work Breakdown}
\label{sec:workbreakdown}
\begin{table}[H]
\centering
\caption{Milestone task allocation (see also the project catalog entry for this
project).}
\label{tab:tasks}
\begin{tabular}{@{}p{0.12\linewidth}p{0.26\linewidth}p{0.26\linewidth}p{0.26\linewidth}@{}}
\toprule
\textbf{Milestone} & \textbf{Student A} & \textbf{Student B} & \textbf{Student C} \\
\midrule
M1 (Wk 3) & JPEG-DCT extraction pipeline; dataset subset chosen & Design DCT-domain classifier architecture & Define accuracy/speed comparison protocol \\
M2 (Wk 6) & RGB-ResNet baseline trained, accuracy recorded & DCT-domain classifier, first trained checkpoint & Preliminary accuracy + speed comparison \\
M3 (Wk 9) & Final RGB baseline (accuracy + decode+infer time) & Final DCT-domain classifier & Full accuracy-vs-speed table, ablation on which DCT bands matter \\
\bottomrule
\end{tabular}
\end{table}

\bibliographystyle{unsrtnat}
\bibliography{references}

\end{document}
